Indeed, (ii) implies that there exists an effective divisor @@M0@@ such that @@M1@@, and @@M2@@ implies that @@M3@@; whence (i).

To see that (i') implies (ii), we may clearly suppose that @@M4@@ is very ample; then there exists a closed integral, locally principal curve @@M5@@ such that @@M6@@. By Riemann's theorem, (i') implies that @@M7@@ for @@M8@@ and @@M9@@ with @@M10@@. Since @@M11@@ for @@M12@@, the exact sequence
@@M13@@
gives @@M14@@ for @@M15@@. Thus, for @@M16@@, we have
@@M17@@
which completes the proof of the lemma.

\sgaStableTarget{sga6.e13.para.7-1-3}\textbf{7.1.3.} Now suppose in 7.1 that there exists an invertible sheaf @@M18@@ on @@M19@@ such that @@M20@@, and argue by contradiction. Replace @@M21@@ by the underlying reduced scheme of a Chow presentation of @@M22@@, which is allowed by the projection formula X 4.4.4, and let @@M23@@ be ample on @@M24@@. Take @@M25@@ such that @@M26@@ is ample, and such that @@M27@@; take @@M28@@ (@@M29@@) such that, putting @@M30@@, we have that @@M31@@ is @@M32@@; then we also have that @@M33@@ is @@M34@@; this returns us to case 7.1.1, with @@M35@@ and @@M36@@.

\sgaStableTarget{sga6.e13.para.7-1-4}\textbf{7.1.4.} Finally, consider the general case of 7.1. Still arguing by contradiction, suppose that @@M37@@ is not numerically equivalent to @@M38@@; in addition, replace @@M39@@ by the normalization of @@M40@@, so that @@M41@@ has
